\section{Experimental evidence register}\label{sec:tt-evidence}
The next step is a primary-literature search and evidence audit. The table is deliberately empty of claimed measurements until a source and its experimental conditions have been checked.

\begin{table}[htbp]
\centering
\caption{Evidence extraction schema. ``Pending'' means \emph{not yet verified}, not that measurements do not exist.}\label{tab:tt-evidence}
\small
\begin{tabular}{p{0.17\linewidth}p{0.27\linewidth}p{0.26\linewidth}p{0.19\linewidth}}
\toprule
Source & Observable / kinematics & Reported extraction & Status\\
\midrule
Primary literature audit & $\dd\sigma/\dd\Omega$, polarization, or other published quantity & Phase shifts, $a_C$, $r_C$, or model constraints if actually reported & Pending\\
\bottomrule
\end{tabular}
\end{table}

\subsection{Extraction protocol}
For each candidate paper record the apparatus and target, absolute or relative normalization, angular coverage, beam-energy definition, quoted systematic errors, Coulomb subtraction, fitted partial waves, and whether the result is an experimental datum or a reanalysis of someone else's data. In a final synthesis, estimates without a reported uncertainty should be separated from constrained intervals.

\subsection{Identifiability}
A sparse set of differential cross sections may constrain only combinations of partial-wave amplitudes. Fit covariances, discrete ambiguities, resolution smearing, and the choice of reaction channels can dominate the uncertainty of threshold extrapolations. Such uncertainties must not disappear when translating a model fit into a single quoted ``potential parameter.''
